\begin{abstract}
Se comparó una red convolucional de referencia publicada en Kaggle con YOLO26n-cls para clasificar imágenes renales en quiste, normal, cálculo y tumor. Se descargó la versión 2 de la recopilación de GonzaJL y se auditaron sus archivos antes de formar las particiones. De 11.756 imágenes legibles se retiraron 1.225 copias exactas adicionales. Las 10.531 imágenes restantes se distribuyeron en 7.372 de entrenamiento, 2.633 de prueba y 526 de validación, manteniendo los grupos de similitud dentro de un único conjunto. Cada arquitectura se entrenó desde cero con cinco semillas, imágenes de 224 por 224, Adam, un máximo de 20 épocas y paciencia de cinco. La selección del checkpoint se realizó mediante F1 macro de validación. En prueba, la CNN obtuvo F1 macro medio 0{,}7077 y YOLO26n-cls 0{,}7478. La diferencia media YOLO menos CNN fue 0{,}0400, con valor p bilateral exacto 0{,}1875. La latencia mediana media fue 1{,}5233 ms para la CNN y 6{,}3146 ms para YOLO. Las cinco repeticiones describen la variabilidad del entrenamiento sobre una partición fija. La interpretación del desempeño queda limitada por la procedencia heterogénea y la ausencia de identificadores verificables de paciente.
\end{abstract}
